% ════════════════════════════════════════════
% 数据与变量处理（5.1.2）
% ════════════════════════════════════════════
\subsubsection{数据与变量处理}

数据理解阶段的目标是如实记录原始状态，任何填补或变换都会掩盖分布问题，因此本部分只做类型转换与质量核查，不改变数值。处理时先以只读方式载入原始文件，保留 500 行、8 列的全部记录，并附加行号作为唯一标识，避免后续排序造成错位。随后按字段含义转换数据类型，把帖子类型、内容主题、星期与时段声明为分类变量，把粉丝规模、曝光次数与发帖小时声明为数值变量。缺失情况按列统计缺失条数占 500 条记录的比例，用以判断问题集中在哪些字段；重复记录按 8 个原始字段逐行比对；取值范围则与字段说明核对，对越界值只标记、不修改，把处理决定留到数据整理阶段统一执行。

经过上述处理，8 个字段的类型与缺失状态全部确认，原始数值未被改写，后续清洗可以在同一份数据上继续。由于本部分不做变换，前后对比图没有实质内容，曝光量的分布形态改由结果部分的图\ref{fig:q1_dist}呈现。

在明确数据状态后，需要选择能够回答“数据结构与变量关系”这一问题的分析工具。针对这一任务，需在描述统计、相关分析、方差分解与可视化等方法中择优，从分析目标、输出形式与适用边界三个维度对比其适用性，如表\ref{tab:algo1}所示。

\begin{longtable}{>{\centering\arraybackslash}p{0.16\textwidth}>{\centering\arraybackslash}p{0.26\textwidth}>{\centering\arraybackslash}p{0.22\textwidth}>{\centering\arraybackslash}p{0.24\textwidth}}
  \caption{数据理解方法对比}
  \label{tab:algo1} \\
  \toprule
  方法 & 主要用途 & 输出形式 & 适用边界 \\
  \midrule
  \endfirsthead
  \caption{数据理解方法对比（续）} \\
  \toprule
  方法 & 主要用途 & 输出形式 & 适用边界 \\
  \midrule
  \endhead
  \bottomrule
  \endfoot
  描述统计 & 刻画集中趋势、离散程度与尾部厚度 & 均值、分位数、偏度、峰度 & 不能说明变量间关系 \\
  相关分析 & 度量两两线性与单调关联 & Pearson、Spearman 系数 & 不能控制其他变量 \\
  方差分解 & 判断变异来自组间还是组内 & 组间与组内方差、解释度 & 只针对单一分组变量 \\
  可视化 & 呈现重尾、分组差异与共线形态 & 直方图、箱线图、散点图 & 偏主观，需数值佐证 \\
\end{longtable}

如表\ref{tab:algo1}所示，四类方法各有所长，单独使用都不足以支撑变量取舍。描述统计能说明曝光量是否重尾，相关分析能说明粉丝规模与月份是否同步，方差分解能进一步回答“月内是否还有可用变异”，可视化则用于检查数值结论是否由少数点驱动。综合以上分析，本部分选择以描述统计与方差分解为主、相关分析与可视化辅助的组合方案。
